\section{Significance of the Study}
\label{sec:significance_of_the_study}

This study contributes to reinforcement-learning research by examining an intrinsic-reward design that explicitly separates surprise from learning progress.
The distinction is significant because high prediction error alone does not necessarily imply useful experience.
By emphasizing local improvement in a dynamics model, the proposed method addresses a major limitation of curiosity-driven exploration: the tendency to reward transitions that are unpredictable but not productive.

The study is also significant for sparse-reward learning.
In environments where successful behavior requires long action sequences before substantial external reward is observed, an agent requires additional signals that make early learning possible.
An intrinsic reward based on both impact and learning progress can guide the agent toward transitions that change the learned representation while still favoring regions where predictive competence is improving.
This provides a principled alternative to using novelty or prediction error alone.

For continuous-control problems, the study demonstrates how region-local exploration statistics can be maintained in a learned latent space rather than in a manually discretized state space.
This is important because continuous observations do not naturally support simple tabular visitation counts or fixed region definitions.
An online partition of feature space allows the method to adapt its locality structure as training data are collected.

The evaluation framework further contributes by considering both sample efficiency and computational efficiency.
Intrinsic rewards can improve learning per environment step while increasing the time required for each update.
Reporting both step-normalized and wall-clock-normalized performance gives a more complete view of whether an exploration method is practically useful.
The inclusion of ablations also clarifies how each component of the proposed method affects learning, reliability, and computational tradeoffs.

Finally, the study documents a complete reinforcement-learning research workflow: theoretical formulation, method design, experimental comparison, ablation analysis, and interpretation of results.
The work is therefore relevant to research on exploration, intrinsic motivation, and empirical evaluation in deep reinforcement learning.